\documentclass[11pt]{article}
\usepackage[T1]{fontenc}
\usepackage[margin=1in]{geometry}
\usepackage{amsmath,amssymb,amsthm,microtype,graphicx,mathtools}
\usepackage[colorlinks=true,urlcolor=blue,citecolor=blue,linkcolor=blue]{hyperref}
\newtheorem{theorem}{Theorem}
\newtheorem{remark}[theorem]{Remark}
\newcommand{\Mthree}{\mathbb M_3}
\newcommand{\Cstar}{C^*}
\DeclareMathOperator{\spanop}{span}
\title{A generating counterexample to the pure extension property in \(\mathbb M_3\)}
\author{Counterexample packet for \texttt{fa\_banach\_001}}
\date{11 August 2026}

\begin{document}
\maketitle

\noindent\textbf{Status.}
Candidate full negative resolution, subject to human review.  An explicit
three-dimensional unital subspace \(M\subset\mathbb M_3\) generates all of
\(\mathbb M_3\), yet a pure state on \(M\) has two distinct pure state
extensions.  The same phenomenon holds for the four-dimensional
operator-system hull \(M+M^*\).

\section*{Source question}

Rapha\"el Clou\^atre, \emph{Restrictions of pure states to subspaces of
\(C^*\)-algebras}, arXiv:2306.02365, \emph{Journal of Functional Analysis}
289 (2025), no.\ 10, Article 111104,
\href{https://doi.org/10.1016/j.jfa.2025.111104}{doi:10.1016/j.jfa.2025.111104}
\cite{clouatre}.  At the end of Section 2.3 (page 10 of the cached revised
source PDF), after proving the result for \(\mathbb M_2\), the paper states:
\begin{quote}
\textquotedblleft In summary, any unital subspace of \(\mathbb M_2\) has the
pure extension property, while there are subspaces of \(\mathbb M_4\) without
this property [57]. We do not know what happens for subspaces in
\(\mathbb M_3\).\textquotedblright
\end{quote}

\begin{center}
\includegraphics[width=0.96\linewidth]{figures/open_problem_crop.png}
\end{center}

\section*{Definitions}

For a unital subspace \(M\) of a unital \(C^*\)-algebra \(B\), a state on
\(M\) is a unital contractive linear functional.  It is \emph{pure} when it
is an extreme point of the state space.  The space \(M\) has the
\emph{pure extension property in \(B\)} if every pure state on \(M\) has a
unique state extension to \(B\) \cite[\S2]{clouatre}.  A state on a unital
subspace extends to a state on the ambient \(C^*\)-algebra: extend it
norm-preservingly by Hahn--Banach, then use the standard fact that a unital
norm-one functional on a \(C^*\)-algebra is positive.

Let \(E_{jk}\) denote the standard matrix units in \(\Mthree\), and let
\(e_1,e_2,e_3\) be the standard basis of \(\mathbb C^3\).

\section*{Proof intuition}

The construction creates a two-dimensional top eigenspace on which two
different vector states cannot be distinguished, while two independent
couplings to the third coordinate force the generated \(C^*\)-algebra to be
all of \(\Mthree\).  The top eigenspace is cut out by
\(H=\operatorname{diag}(1,1,0)\).  Maximizing the expectation of \(H\)
forces every extending density matrix to be supported on that eigenspace.
The coupling matrices then have zero compression there, so all such
extensions induce the same functional on the chosen subspace.  Consequently
that functional is not merely extreme but exposed.  Nevertheless, the
vector states at \(e_1\) and \(e_2\) remain distinct on \(\Mthree\).

\section*{Counterexample}

Set
\[
 H=E_{11}+E_{22},\qquad
 X=E_{13}+E_{31},\qquad
 Y=E_{23}+E_{32},
\]
and put
\[
 A=H+iX,\qquad B=H+iY,\qquad
 M=\spanop_{\mathbb C}\{I,A,B\}\subset\Mthree .
\]

\begin{theorem}\label{thm:main}
The subspace \(M\) has complex dimension \(3\) and
\(\Cstar(M)=\Mthree\).  The common restriction
\[
 \varphi=\omega_1|_M=\omega_2|_M,\qquad
 \omega_j(T)=\langle Te_j,e_j\rangle,
\]
is an exposed, hence pure, state on \(M\), but \(\omega_1\) and \(\omega_2\)
are distinct pure state extensions of \(\varphi\) to \(\Mthree\).
Consequently \(M\) does not have the pure extension property in
\(\Cstar(M)=\Mthree\).

Moreover,
\[
 S=M+M^*=\spanop_{\mathbb C}\{I,H,X,Y\}
\]
is a four-dimensional operator system with \(\Cstar(S)=\Mthree\), and the
common restriction of \(\omega_1,\omega_2\) to \(S\) is exposed and pure.
Thus the failure persists even within the class of generating operator
systems in \(\Mthree\).
\end{theorem}

\begin{proof}
\textbf{Generation.}
The matrices \(I,A,B\) are linearly independent because the \(E_{13}\) and
\(E_{23}\) entries first force the coefficients of \(A\) and \(B\) to
vanish.  Hence \(\dim_{\mathbb C}M=3\).
Because a generated \(C^*\)-algebra contains adjoints,
\[
 H=\frac{A+A^*}{2},\quad X=\frac{A-A^*}{2i},\quad
 Y=\frac{B-B^*}{2i}
\]
belong to \(\Cstar(M)\).  Thus \(E_{33}=I-H\) belongs to \(\Cstar(M)\), and
\[
 XE_{33}=E_{13},\quad E_{33}X=E_{31},\quad
 YE_{33}=E_{23},\quad E_{33}Y=E_{32}.
\]
Products give
\[
 E_{13}E_{32}=E_{12},\qquad E_{23}E_{31}=E_{21},
\]
as well as \(E_{11}=E_{13}E_{31}\) and
\(E_{22}=E_{23}E_{32}\).  All nine matrix units therefore lie in
\(\Cstar(M)\), so \(\Cstar(M)=\Mthree\).

\textbf{The two extensions agree.}
For \(j=1,2\),
\[
 \omega_j(I)=1,\qquad \omega_j(A)=1,\qquad \omega_j(B)=1.
\]
Thus the restrictions of \(\omega_1\) and \(\omega_2\) to \(M\) coincide;
call the common state \(\varphi\).

\textbf{The common restriction is exposed.}
On the state space of \(M\), consider the real affine functional
\[
 F(\psi)=\operatorname{Re}\frac{\psi(A)+\psi(B)}2.
\]
Let \(\psi\) be any state on \(M\), and let \(\rho\) be a state extension to
\(\Mthree\).  Since \(X\) and \(Y\) are self-adjoint,
\[
 F(\psi)=\rho(H)\leq 1.
\]
Suppose equality holds.  Write
\(\rho(T)=\operatorname{Tr}(DT)\) for a density matrix \(D\geq0\) with
\(\operatorname{Tr}D=1\).  Then
\[
 1=\rho(H)=D_{11}+D_{22},\qquad D_{33}=0.
\]
Positivity gives
\(\lvert D_{j3}\rvert^2\leq D_{jj}D_{33}=0\) for \(j=1,2\).  Hence the third
row and column of \(D\) vanish, so
\[
 \rho(X)=\rho(Y)=0.
\]
It follows that \(\psi(A)=\psi(B)=1\), and therefore \(\psi=\varphi\) on the
basis \(I,A,B\).  Thus \(\varphi\) is the unique maximizer of \(F\), so it is
an exposed state and in particular is pure.

\textbf{Nonuniqueness.}
Each \(\omega_j\) is a vector state in the irreducible defining
representation of \(\Mthree\), hence is pure.  They are distinct because
\[
 \omega_1(E_{11})=1,\qquad \omega_2(E_{11})=0.
\]
They are therefore two distinct pure extensions of the pure state
\(\varphi\), proving the failure of the pure extension property.

\textbf{Operator-system strengthening.}
The displayed adjoint formulas show
\[
 M+M^*=\spanop_{\mathbb C}\{I,H,X,Y\}=S.
\]
These four matrices are linearly independent and self-adjoint, and the
generation argument already proves \(\Cstar(S)=\Mthree\).  The restrictions
of \(\omega_1,\omega_2\) to \(S\) agree on the displayed basis.  For a state
\(\sigma\) on \(S\), the affine functional \(\sigma\mapsto\sigma(H)\) is at
most \(1\); the density-matrix argument above shows that equality forces
\(\sigma(I,H,X,Y)=(1,1,0,0)\).  Hence the common restriction is again the
unique maximizer, and is exposed and pure.  Its two extensions remain
distinct.
\end{proof}

\section*{Verification notes}

\begin{itemize}
\item The proof works in the strongest natural ambient interpretation:
      both \(M\) and \(S\) generate all of \(\Mthree\), so the ambient
      algebra equals \(\Cstar(M)=\Cstar(S)\).
\item No numerical approximation is used.  The optional checker supplied in
      the packet's code directory verifies Hermiticity, the two vector-state
      restrictions, and that the generated matrix-unit span has complex
      dimension \(9\).
\item The key purity step does not infer purity from purity of an extension.
      It proves directly that the common restriction is the unique maximizer
      of a real affine functional on the state space.
\item Positivity is used only through the elementary density-matrix fact
      \(D_{33}=0\Rightarrow D_{13}=D_{23}=0\).
\end{itemize}

\section*{Novelty check and limitations}

The cached source is the 31 October 2024 advanced revision.  A bounded search
performed on 11 August 2026 checked the arXiv record, the 2025 published
version, and exact phrase/keyword combinations involving pure extension
property, \(\mathbb M_3\), operator systems, and the paper title.
The published 2025 version still states the \(\mathbb M_3\) question as open,
and no later exact resolution was found.  Searches also inspected the
authors' later operator-system work returned by these queries, but found no
answer to this finite-dimensional question.  This supports only a
candidate-novel status; a specialist citation search remains appropriate
before publication.

The theorem resolves the precise finite-dimensional question negatively.  It
does not address the paper's separate questions about the full
noncommutative disc or Drury--Arveson algebras, nor its broader classification
of uniquely extendible restrictions.

\section*{Human-review recommendation}

High confidence in validity: the construction and proof reduce to explicit
\(3\times3\) matrix identities and a one-line positive-semidefinite argument.
The primary review focus should be terminological---confirm that the paper's
state convention for a nonself-adjoint unital subspace is the standard
unital-contractive one used above---and novelty, because the result is
surprisingly elementary.  The operator-system strengthening independently
avoids that convention and should be checked as the most robust form.

\begin{thebibliography}{9}
\bibitem{clouatre}
R. Clou\^atre,
\emph{Restrictions of pure states to subspaces of \(C^*\)-algebras},
J. Funct. Anal. \textbf{289} (2025), no.~10, 111104;
\href{https://arxiv.org/abs/2306.02365}{arXiv:2306.02365}.
\end{thebibliography}

\end{document}
